\documentclass[12pt]{article}
\usepackage[utf8]{inputenc}
\usepackage[T1]{fontenc}
\usepackage{hyperref}
\title{Kenaboi Language Resources for K-12 Curriculum Development}
\author{Compiled by Claude Sonnet 5.5 (claude-sonnet-5-5) for Andr\'as Kornai}
\date{}
\begin{document}
\maketitle

\section{Introduction}
Kenaboi is an extinct language of Negeri Sembilan, Malaysia. It is unclassified or an isolate. It is known only from word lists of about 250 items, collected around 1870--1890 by D.F.A. Hervey and published in Skeat and Blagden's \emph{Pagan Races of the Malay Peninsula} (1906). Glottolog treats it as an isolate. Wikipedia says it may instead be an Aslian (Austroasiatic) language. The speakers appear to have merged into the Temuan Orang Asli community around 1880. There are no native speakers and no known teaching materials. A curriculum would have to be a heritage-awareness programme built on the small surviving word list.

\section{Primary Educational Resources}
I found no primers, textbooks or school programmes for Kenaboi in search results.

\section{Linguistic and Reference Materials}
\begin{itemize}
\item Wikipedia, ``Kenaboi language'': overview of the data and the classification debate. \url{https://en.wikipedia.org/wiki/Kenaboi_language}
\item Glottolog entry for Kenaboi (kena1236), with classification notes. \url{https://glottolog.org/resource/languoid/id/kena1236}
\item Austronesian Basic Vocabulary Database, Kenaboi (I) and Kenaboi (II) word lists. \url{https://abvd.eva.mpg.de/austronesian/language.php?id=1564} and \url{https://abvd.eva.mpg.de/austronesian/language.php?id=1568}
\item The main scholarly source is John Hajek (1998), ``Kenaboi: An extinct unclassified language of the Malay Peninsula'', \emph{Mon-Khmer Studies} 28. It was mentioned in the search summary, but I saw no direct URL for it.
\end{itemize}

\section{Religious and Community Resources}
I found no Kenaboi community organisation or religious materials. Free Malaysia Today notes that the speakers were absorbed into the Temuan community. \url{https://www.freemalaysiatoday.com/category/2022/02/21/languages-that-either-went-extinct-survived-or-were-revived}

\section{Digital and Audio Resources}
I found no audio recordings or software for Kenaboi. No one knows how the words were pronounced.

\section{Curriculum Development Context}
Very little exists. A K-12 unit could cover the language's history, the Orang Asli context, how word lists are used to study extinct languages, and the ~250 recorded words. Teaching Kenaboi as a spoken language would not be honest, because pronunciation, grammar and texts are unknown. Any such work should involve the Temuan and other Orang Asli communities.
\end{document}